\chapter{Migration Guide: Bijou v4.0.0 To v4.1.0}\label{record-264}

\noindent{\footnotesize\color{BookMuted}\path{docs/releases/4.1.0/migration-guide.md}\par}\smallskip

This guide is for users already on \texttt{v\allowbreak{}4\allowbreak{}.\allowbreak{}0\allowbreak{}.\allowbreak{}0} who want to move to \texttt{v\allowbreak{}4\allowbreak{}.\allowbreak{}1\allowbreak{}.\allowbreak{}0}.

If you are coming from an older major line, read \href{https://github.com/flyingrobots/bijou/blob/cbd70612140e7a450d764b3c7be1c6af878b99bc/docs/MIGRATING_TO_V4.md}{\texttt{M\allowbreak{}I\allowbreak{}G\allowbreak{}R\allowbreak{}A\allowbreak{}T\allowbreak{}I\allowbreak{}N\allowbreak{}G\allowbreak{}\_\allowbreak{}T\allowbreak{}O\allowbreak{}\_\allowbreak{}V\allowbreak{}4\allowbreak{}.\allowbreak{}m\allowbreak{}d}} first.

\section{Migration At A Glance}\label{record-264-migration-at-a-glance}

Not every user has real migration work to do.

{\def\LTcaptype{none} % do not increment counter
\begin{longtable}[]{@{}
  >{\raggedright\arraybackslash}p{(\linewidth - 4\tabcolsep) * \real{0.3333}}
  >{\raggedright\arraybackslash}p{(\linewidth - 4\tabcolsep) * \real{0.3333}}
  >{\raggedright\arraybackslash}p{(\linewidth - 4\tabcolsep) * \real{0.3333}}@{}}
\toprule\noalign{}
{I\allowbreak{}f} {y\allowbreak{}o\allowbreak{}u} {u\allowbreak{}s\allowbreak{}e\allowbreak{}…}
 & {E\allowbreak{}x\allowbreak{}p\allowbreak{}e\allowbreak{}c\allowbreak{}t\allowbreak{}e\allowbreak{}d} {m\allowbreak{}i\allowbreak{}g\allowbreak{}r\allowbreak{}a\allowbreak{}t\allowbreak{}i\allowbreak{}o\allowbreak{}n} {e\allowbreak{}f\allowbreak{}f\allowbreak{}o\allowbreak{}r\allowbreak{}t}
 & {W\allowbreak{}h\allowbreak{}y}
 \\
\midrule\noalign{}
\endhead
\bottomrule\noalign{}
\endlastfoot
\texttt{@\allowbreak{}f\allowbreak{}l\allowbreak{}y\allowbreak{}i\allowbreak{}n\allowbreak{}g\allowbreak{}r\allowbreak{}o\allowbreak{}b\allowbreak{}o\allowbreak{}t\allowbreak{}s\allowbreak{}/\allowbreak{}b\allowbreak{}i\allowbreak{}j\allowbreak{}o\allowbreak{}u} {f\allowbreak{}o\allowbreak{}r} {C\allowbreak{}L\allowbreak{}I} {o\allowbreak{}u\allowbreak{}t\allowbreak{}p\allowbreak{}u\allowbreak{}t\allowbreak{},} {p\allowbreak{}r\allowbreak{}o\allowbreak{}m\allowbreak{}p\allowbreak{}t\allowbreak{}s\allowbreak{},} {a\allowbreak{}n\allowbreak{}d} {s\allowbreak{}t\allowbreak{}r\allowbreak{}i\allowbreak{}n\allowbreak{}g\allowbreak{}-\allowbreak{}f\allowbreak{}i\allowbreak{}r\allowbreak{}s\allowbreak{}t} {f\allowbreak{}l\allowbreak{}o\allowbreak{}w\allowbreak{}s} & {L\allowbreak{}o\allowbreak{}w} & {M\allowbreak{}o\allowbreak{}s\allowbreak{}t} \texttt{v\allowbreak{}4\allowbreak{}.\allowbreak{}1\allowbreak{}.\allowbreak{}0} {w\allowbreak{}o\allowbreak{}r\allowbreak{}k} {i\allowbreak{}s} {a\allowbreak{}d\allowbreak{}d\allowbreak{}i\allowbreak{}t\allowbreak{}i\allowbreak{}v\allowbreak{}e} {o\allowbreak{}r} {d\allowbreak{}o\allowbreak{}c\allowbreak{}s\allowbreak{}-\allowbreak{}f\allowbreak{}a\allowbreak{}c\allowbreak{}i\allowbreak{}n\allowbreak{}g} {f\allowbreak{}o\allowbreak{}r} {t\allowbreak{}h\allowbreak{}i\allowbreak{}s} {s\allowbreak{}u\allowbreak{}r\allowbreak{}f\allowbreak{}a\allowbreak{}c\allowbreak{}e\allowbreak{}.} \\
\texttt{@\allowbreak{}f\allowbreak{}l\allowbreak{}y\allowbreak{}i\allowbreak{}n\allowbreak{}g\allowbreak{}r\allowbreak{}o\allowbreak{}b\allowbreak{}o\allowbreak{}t\allowbreak{}s\allowbreak{}/\allowbreak{}b\allowbreak{}i\allowbreak{}j\allowbreak{}o\allowbreak{}u\allowbreak{}-\allowbreak{}t\allowbreak{}u\allowbreak{}i} {f\allowbreak{}o\allowbreak{}r} {f\allowbreak{}u\allowbreak{}l\allowbreak{}l\allowbreak{}s\allowbreak{}c\allowbreak{}r\allowbreak{}e\allowbreak{}e\allowbreak{}n} {r\allowbreak{}u\allowbreak{}n\allowbreak{}t\allowbreak{}i\allowbreak{}m\allowbreak{}e} {a\allowbreak{}p\allowbreak{}p\allowbreak{}s} & {M\allowbreak{}e\allowbreak{}d\allowbreak{}i\allowbreak{}u\allowbreak{}m} & {T\allowbreak{}h\allowbreak{}e} {r\allowbreak{}u\allowbreak{}n\allowbreak{}t\allowbreak{}i\allowbreak{}m\allowbreak{}e} {a\allowbreak{}n\allowbreak{}d} {c\allowbreak{}o\allowbreak{}m\allowbreak{}m\allowbreak{}a\allowbreak{}n\allowbreak{}d} {m\allowbreak{}o\allowbreak{}d\allowbreak{}e\allowbreak{}l} {a\allowbreak{}r\allowbreak{}e} {m\allowbreak{}o\allowbreak{}r\allowbreak{}e} {e\allowbreak{}x\allowbreak{}p\allowbreak{}l\allowbreak{}i\allowbreak{}c\allowbreak{}i\allowbreak{}t\allowbreak{},} {a\allowbreak{}n\allowbreak{}d} {y\allowbreak{}o\allowbreak{}u} {s\allowbreak{}h\allowbreak{}o\allowbreak{}u\allowbreak{}l\allowbreak{}d} {a\allowbreak{}l\allowbreak{}i\allowbreak{}g\allowbreak{}n} {w\allowbreak{}i\allowbreak{}t\allowbreak{}h} {t\allowbreak{}h\allowbreak{}e} {n\allowbreak{}e\allowbreak{}w\allowbreak{}e\allowbreak{}r} {s\allowbreak{}e\allowbreak{}a\allowbreak{}m\allowbreak{}s\allowbreak{}.} \\
\texttt{@\allowbreak{}f\allowbreak{}l\allowbreak{}y\allowbreak{}i\allowbreak{}n\allowbreak{}g\allowbreak{}r\allowbreak{}o\allowbreak{}b\allowbreak{}o\allowbreak{}t\allowbreak{}s\allowbreak{}/\allowbreak{}b\allowbreak{}i\allowbreak{}j\allowbreak{}o\allowbreak{}u\allowbreak{}-\allowbreak{}t\allowbreak{}u\allowbreak{}i\allowbreak{}-\allowbreak{}a\allowbreak{}p\allowbreak{}p} {/} \texttt{c\allowbreak{}r\allowbreak{}e\allowbreak{}a\allowbreak{}t\allowbreak{}e\allowbreak{}F\allowbreak{}r\allowbreak{}a\allowbreak{}m\allowbreak{}e\allowbreak{}d\allowbreak{}A\allowbreak{}p\allowbreak{}p\allowbreak{}(\allowbreak{})} {h\allowbreak{}e\allowbreak{}a\allowbreak{}v\allowbreak{}i\allowbreak{}l\allowbreak{}y} & {M\allowbreak{}e\allowbreak{}d\allowbreak{}i\allowbreak{}u\allowbreak{}m} {t\allowbreak{}o} {H\allowbreak{}i\allowbreak{}g\allowbreak{}h} & {F\allowbreak{}r\allowbreak{}a\allowbreak{}m\allowbreak{}e\allowbreak{}d\allowbreak{}-\allowbreak{}s\allowbreak{}h\allowbreak{}e\allowbreak{}l\allowbreak{}l} {t\allowbreak{}y\allowbreak{}p\allowbreak{}i\allowbreak{}n\allowbreak{}g\allowbreak{},} {l\allowbreak{}a\allowbreak{}y\allowbreak{}e\allowbreak{}r} {r\allowbreak{}o\allowbreak{}u\allowbreak{}t\allowbreak{}i\allowbreak{}n\allowbreak{}g\allowbreak{},} {a\allowbreak{}n\allowbreak{}d} {s\allowbreak{}h\allowbreak{}e\allowbreak{}l\allowbreak{}l} {i\allowbreak{}n\allowbreak{}s\allowbreak{}p\allowbreak{}e\allowbreak{}c\allowbreak{}t\allowbreak{}i\allowbreak{}o\allowbreak{}n} {s\allowbreak{}u\allowbreak{}r\allowbreak{}f\allowbreak{}a\allowbreak{}c\allowbreak{}e\allowbreak{}s} {a\allowbreak{}r\allowbreak{}e} {m\allowbreak{}o\allowbreak{}r\allowbreak{}e} {e\allowbreak{}x\allowbreak{}p\allowbreak{}l\allowbreak{}i\allowbreak{}c\allowbreak{}i\allowbreak{}t\allowbreak{},} {s\allowbreak{}o} {o\allowbreak{}l\allowbreak{}d} {a\allowbreak{}s\allowbreak{}s\allowbreak{}u\allowbreak{}m\allowbreak{}p\allowbreak{}t\allowbreak{}i\allowbreak{}o\allowbreak{}n\allowbreak{}s} {m\allowbreak{}a\allowbreak{}y} {n\allowbreak{}e\allowbreak{}e\allowbreak{}d} {c\allowbreak{}l\allowbreak{}e\allowbreak{}a\allowbreak{}n\allowbreak{}u\allowbreak{}p\allowbreak{}.} \\
{c\allowbreak{}u\allowbreak{}s\allowbreak{}t\allowbreak{}o\allowbreak{}m} {r\allowbreak{}u\allowbreak{}n\allowbreak{}t\allowbreak{}i\allowbreak{}m\allowbreak{}e} {i\allowbreak{}n\allowbreak{}t\allowbreak{}e\allowbreak{}r\allowbreak{}n\allowbreak{}a\allowbreak{}l\allowbreak{}s} {o\allowbreak{}r} {s\allowbreak{}h\allowbreak{}e\allowbreak{}l\allowbreak{}l\allowbreak{}-\allowbreak{}a\allowbreak{}d\allowbreak{}j\allowbreak{}a\allowbreak{}c\allowbreak{}e\allowbreak{}n\allowbreak{}t} {h\allowbreak{}e\allowbreak{}l\allowbreak{}p\allowbreak{}e\allowbreak{}r\allowbreak{}s} & {H\allowbreak{}i\allowbreak{}g\allowbreak{}h} & {T\allowbreak{}h\allowbreak{}e} {r\allowbreak{}e\allowbreak{}p\allowbreak{}o} {i\allowbreak{}s} {m\allowbreak{}o\allowbreak{}v\allowbreak{}i\allowbreak{}n\allowbreak{}g} {t\allowbreak{}o\allowbreak{}w\allowbreak{}a\allowbreak{}r\allowbreak{}d} {r\allowbreak{}u\allowbreak{}n\allowbreak{}t\allowbreak{}i\allowbreak{}m\allowbreak{}e\allowbreak{}-\allowbreak{}e\allowbreak{}n\allowbreak{}g\allowbreak{}i\allowbreak{}n\allowbreak{}e} {s\allowbreak{}e\allowbreak{}a\allowbreak{}m\allowbreak{}s} {a\allowbreak{}n\allowbreak{}d} {a\allowbreak{}w\allowbreak{}a\allowbreak{}y} {f\allowbreak{}r\allowbreak{}o\allowbreak{}m} {b\allowbreak{}r\allowbreak{}a\allowbreak{}n\allowbreak{}c\allowbreak{}h\allowbreak{}-\allowbreak{}o\allowbreak{}r\allowbreak{}d\allowbreak{}e\allowbreak{}r} {s\allowbreak{}h\allowbreak{}e\allowbreak{}l\allowbreak{}l} {f\allowbreak{}o\allowbreak{}l\allowbreak{}k\allowbreak{}l\allowbreak{}o\allowbreak{}r\allowbreak{}e\allowbreak{}.} \\
{l\allowbreak{}o\allowbreak{}c\allowbreak{}a\allowbreak{}l\allowbreak{}i\allowbreak{}z\allowbreak{}a\allowbreak{}t\allowbreak{}i\allowbreak{}o\allowbreak{}n} {w\allowbreak{}o\allowbreak{}r\allowbreak{}k\allowbreak{}f\allowbreak{}l\allowbreak{}o\allowbreak{}w\allowbreak{}s} & {O\allowbreak{}p\allowbreak{}t\allowbreak{}i\allowbreak{}o\allowbreak{}n\allowbreak{}a\allowbreak{}l} {b\allowbreak{}u\allowbreak{}t} {m\allowbreak{}e\allowbreak{}a\allowbreak{}n\allowbreak{}i\allowbreak{}n\allowbreak{}g\allowbreak{}f\allowbreak{}u\allowbreak{}l} & {E\allowbreak{}n\allowbreak{}t\allowbreak{}i\allowbreak{}r\allowbreak{}e\allowbreak{}l\allowbreak{}y} {n\allowbreak{}e\allowbreak{}w} {p\allowbreak{}a\allowbreak{}c\allowbreak{}k\allowbreak{}a\allowbreak{}g\allowbreak{}e\allowbreak{}s} {a\allowbreak{}n\allowbreak{}d} {w\allowbreak{}o\allowbreak{}r\allowbreak{}k\allowbreak{}f\allowbreak{}l\allowbreak{}o\allowbreak{}w\allowbreak{}s} {e\allowbreak{}x\allowbreak{}i\allowbreak{}s\allowbreak{}t} {n\allowbreak{}o\allowbreak{}w} {i\allowbreak{}f} {y\allowbreak{}o\allowbreak{}u} {w\allowbreak{}a\allowbreak{}n\allowbreak{}t} {t\allowbreak{}h\allowbreak{}e\allowbreak{}m\allowbreak{}.} \\
\end{longtable}
}

\section{Upgrade Checklist}\label{record-264-upgrade-checklist}

\begin{enumerate}
\def\labelenumi{\arabic{enumi}.}
\tightlist
\item
  Upgrade all Bijou packages together to the same version.
\item
  Rebuild and re-typecheck before making behavioral changes.
\item
  Audit framed-shell typing if you use \texttt{c\allowbreak{}r\allowbreak{}e\allowbreak{}a\allowbreak{}t\allowbreak{}e\allowbreak{}F\allowbreak{}r\allowbreak{}a\allowbreak{}m\allowbreak{}e\allowbreak{}d\allowbreak{}A\allowbreak{}p\allowbreak{}p\allowbreak{}(\allowbreak{})}.
\item
  Audit custom command code if you rely on long-lived effects or subscriptions.
\item
  Replace shell-folklore assumptions with exported layer/runtime helpers where possible.
\item
  Adopt the new localization packages only if you need that surface.
\item
  Re-run your fullscreen and shell interaction tests, especially for help/settings/notification/quit flows.
\end{enumerate}

\section{1. Upgrade In Lock-Step}\label{record-264-upgrade-in-lock-step}

Bijou still versions its published packages in lock-step.

Upgrade the packages you use together:

\begin{CodeBlock}
npm install @flyingrobots/bijou@4.1.0 @flyingrobots/bijou-node@4.1.0
\end{CodeBlock}

Add the fullscreen packages if your app uses them:

\begin{CodeBlock}
npm install @flyingrobots/bijou-tui@4.1.0 @flyingrobots/bijou-tui-app@4.1.0
\end{CodeBlock}

If you adopt the new localization surface, also add the relevant i18n packages:

\begin{CodeBlock}
npm install @flyingrobots/bijou-i18n@4.1.0 @flyingrobots/bijou-i18n-tools@4.1.0
\end{CodeBlock}

If you also need filesystem or XLSX exchange adapters, install those in the same lock-step release:

\begin{CodeBlock}
npm install @flyingrobots/bijou-i18n-tools-node@4.1.0 @flyingrobots/bijou-i18n-tools-xlsx@4.1.0
\end{CodeBlock}

These i18n packages are part of the \texttt{4\allowbreak{}.\allowbreak{}1\allowbreak{}.\allowbreak{}0} automated publish surface, so they land together with the rest of the workspace in this release.

\section{2. Framed Apps: Move Onto The Typed Shell Surface}\label{record-264-framed-apps-move-onto-the-typed-shell-surface}

If your app uses \texttt{c\allowbreak{}r\allowbreak{}e\allowbreak{}a\allowbreak{}t\allowbreak{}e\allowbreak{}F\allowbreak{}r\allowbreak{}a\allowbreak{}m\allowbreak{}e\allowbreak{}d\allowbreak{}A\allowbreak{}p\allowbreak{}p\allowbreak{}(\allowbreak{})}, this is the migration area most likely to matter.

The shell now exports explicit framed message and update helpers:

\begin{itemize}
\tightlist
\item
  \texttt{F\allowbreak{}r\allowbreak{}a\allowbreak{}m\allowbreak{}e\allowbreak{}P\allowbreak{}a\allowbreak{}g\allowbreak{}e\allowbreak{}M\allowbreak{}s\allowbreak{}g\allowbreak{}<\allowbreak{}M\allowbreak{}s\allowbreak{}g\allowbreak{}>}
\item
  \texttt{F\allowbreak{}r\allowbreak{}a\allowbreak{}m\allowbreak{}e\allowbreak{}P\allowbreak{}a\allowbreak{}g\allowbreak{}e\allowbreak{}U\allowbreak{}p\allowbreak{}d\allowbreak{}a\allowbreak{}t\allowbreak{}e\allowbreak{}R\allowbreak{}e\allowbreak{}s\allowbreak{}u\allowbreak{}l\allowbreak{}t\allowbreak{}<\allowbreak{}P\allowbreak{}a\allowbreak{}g\allowbreak{}e\allowbreak{}M\allowbreak{}o\allowbreak{}d\allowbreak{}e\allowbreak{}l\allowbreak{},\allowbreak{} \allowbreak{}M\allowbreak{}s\allowbreak{}g\allowbreak{}>}
\item
  \texttt{F\allowbreak{}r\allowbreak{}a\allowbreak{}m\allowbreak{}e\allowbreak{}d\allowbreak{}A\allowbreak{}p\allowbreak{}p\allowbreak{}M\allowbreak{}s\allowbreak{}g\allowbreak{}<\allowbreak{}M\allowbreak{}s\allowbreak{}g\allowbreak{}>}
\item
  \texttt{F\allowbreak{}r\allowbreak{}a\allowbreak{}m\allowbreak{}e\allowbreak{}d\allowbreak{}A\allowbreak{}p\allowbreak{}p\allowbreak{}U\allowbreak{}p\allowbreak{}d\allowbreak{}a\allowbreak{}t\allowbreak{}e\allowbreak{}R\allowbreak{}e\allowbreak{}s\allowbreak{}u\allowbreak{}l\allowbreak{}t\allowbreak{}<\allowbreak{}P\allowbreak{}a\allowbreak{}g\allowbreak{}e\allowbreak{}M\allowbreak{}o\allowbreak{}d\allowbreak{}e\allowbreak{}l\allowbreak{},\allowbreak{} \allowbreak{}M\allowbreak{}s\allowbreak{}g\allowbreak{}>}
\item
  \texttt{F\allowbreak{}r\allowbreak{}a\allowbreak{}m\allowbreak{}e\allowbreak{}d\allowbreak{}A\allowbreak{}p\allowbreak{}p\allowbreak{}<\allowbreak{}P\allowbreak{}a\allowbreak{}g\allowbreak{}e\allowbreak{}M\allowbreak{}o\allowbreak{}d\allowbreak{}e\allowbreak{}l\allowbreak{},\allowbreak{} \allowbreak{}M\allowbreak{}s\allowbreak{}g\allowbreak{}>}
\end{itemize}

\subsection{What to change}\label{record-264-what-to-change}

\begin{itemize}
\tightlist
\item
  stop treating framed page updates as effectively \texttt{a\allowbreak{}n\allowbreak{}y}
\item
  update explicit type annotations that assumed page messages were only your raw app \texttt{M\allowbreak{}s\allowbreak{}g}
\item
  update tests that asserted the framed shell collapsed everything back to plain \texttt{M\allowbreak{}s\allowbreak{}g}
\end{itemize}

\subsection{Typical before/after mindset}\label{record-264-typical-beforeafter-mindset}

Before:

\begin{itemize}
\tightlist
\item
  the page/shell seam was often typed loosely
\item
  external casts papered over framed-shell complexity
\end{itemize}

After:

\begin{itemize}
\tightlist
\item
  the page/shell seam should stay typed explicitly
\item
  page updates should accept \texttt{F\allowbreak{}r\allowbreak{}a\allowbreak{}m\allowbreak{}e\allowbreak{}P\allowbreak{}a\allowbreak{}g\allowbreak{}e\allowbreak{}M\allowbreak{}s\allowbreak{}g\allowbreak{}<\allowbreak{}M\allowbreak{}s\allowbreak{}g\allowbreak{}>} where appropriate
\item
  framed app updates should be allowed to return \texttt{C\allowbreak{}m\allowbreak{}d\allowbreak{}<\allowbreak{}F\allowbreak{}r\allowbreak{}a\allowbreak{}m\allowbreak{}e\allowbreak{}d\allowbreak{}A\allowbreak{}p\allowbreak{}p\allowbreak{}M\allowbreak{}s\allowbreak{}g\allowbreak{}<\allowbreak{}M\allowbreak{}s\allowbreak{}g\allowbreak{}>\allowbreak{}>\allowbreak{}[\allowbreak{}]}
\end{itemize}

If your code was already mostly inference-driven and you were not fighting the shell types directly, you may have little to change. If you have custom wrapper helpers around \texttt{c\allowbreak{}r\allowbreak{}e\allowbreak{}a\allowbreak{}t\allowbreak{}e\allowbreak{}F\allowbreak{}r\allowbreak{}a\allowbreak{}m\allowbreak{}e\allowbreak{}d\allowbreak{}A\allowbreak{}p\allowbreak{}p\allowbreak{}(\allowbreak{})}, inspect them carefully.

\section{3. Commands: Stop Pretending They Only End In Messages}\label{record-264-commands-stop-pretending-they-only-end-in-messages}

\texttt{C\allowbreak{}m\allowbreak{}d\allowbreak{}<\allowbreak{}M\allowbreak{}>} now models the runtime more honestly.

A command may now:

\begin{itemize}
\tightlist
\item
  complete synchronously or asynchronously
\item
  emit intermediate messages through \texttt{e\allowbreak{}m\allowbreak{}i\allowbreak{}t\allowbreak{}(\allowbreak{}.\allowbreak{}.\allowbreak{}.\allowbreak{})}
\item
  resolve to a final message
\item
  resolve to \texttt{Q\allowbreak{}U\allowbreak{}I\allowbreak{}T}
\item
  return a cleanup handle or cleanup function
\item
  resolve to \texttt{v\allowbreak{}o\allowbreak{}i\allowbreak{}d}
\end{itemize}

\subsection{What to look for}\label{record-264-what-to-look-for}

\begin{itemize}
\tightlist
\item
  custom command wrappers that assume every command returns a final message
\item
  tests that stub commands as only \texttt{P\allowbreak{}r\allowbreak{}o\allowbreak{}m\allowbreak{}i\allowbreak{}s\allowbreak{}e\allowbreak{}<\allowbreak{}M\allowbreak{}s\allowbreak{}g\allowbreak{} \allowbreak{}|\allowbreak{} \allowbreak{}u\allowbreak{}n\allowbreak{}d\allowbreak{}e\allowbreak{}f\allowbreak{}i\allowbreak{}n\allowbreak{}e\allowbreak{}d\allowbreak{} \allowbreak{}|\allowbreak{} \allowbreak{}Q\allowbreak{}U\allowbreak{}I\allowbreak{}T\allowbreak{}>}
\item
  child-app command mapping that accidentally discards cleanup results
\end{itemize}

\subsection{What to do}\label{record-264-what-to-do}

\begin{itemize}
\tightlist
\item
  widen your local command typings if you mirrored the old narrower contract
\item
  prefer the built-in sub-app lifecycle helpers if you were manually coercing command results
\item
  treat cleanup-producing commands as first-class runtime work, not as hacks that need \texttt{a\allowbreak{}n\allowbreak{}y}
\end{itemize}

If you only use built-in helpers like \texttt{t\allowbreak{}i\allowbreak{}c\allowbreak{}k\allowbreak{}(\allowbreak{})}, \texttt{a\allowbreak{}n\allowbreak{}i\allowbreak{}m\allowbreak{}a\allowbreak{}t\allowbreak{}e\allowbreak{}(\allowbreak{})}, or normal message-returning commands, this may be invisible to you. The migration is most relevant for custom runtime work and command composition helpers.

\section{4. Shell Layers: Use The Exported Truth}\label{record-264-shell-layers-use-the-exported-truth}

If you wrote code or tests that reason about shell layering indirectly, update them to prefer the exported shell inspection surface:

\begin{itemize}
\tightlist
\item
  \texttt{d\allowbreak{}e\allowbreak{}s\allowbreak{}c\allowbreak{}r\allowbreak{}i\allowbreak{}b\allowbreak{}e\allowbreak{}F\allowbreak{}r\allowbreak{}a\allowbreak{}m\allowbreak{}e\allowbreak{}L\allowbreak{}a\allowbreak{}y\allowbreak{}e\allowbreak{}r\allowbreak{}S\allowbreak{}t\allowbreak{}a\allowbreak{}c\allowbreak{}k\allowbreak{}(\allowbreak{}.\allowbreak{}.\allowbreak{}.\allowbreak{})}
\item
  \texttt{a\allowbreak{}c\allowbreak{}t\allowbreak{}i\allowbreak{}v\allowbreak{}e\allowbreak{}F\allowbreak{}r\allowbreak{}a\allowbreak{}m\allowbreak{}e\allowbreak{}L\allowbreak{}a\allowbreak{}y\allowbreak{}e\allowbreak{}r\allowbreak{}(\allowbreak{}.\allowbreak{}.\allowbreak{}.\allowbreak{})}
\item
  \texttt{u\allowbreak{}n\allowbreak{}d\allowbreak{}e\allowbreak{}r\allowbreak{}l\allowbreak{}y\allowbreak{}i\allowbreak{}n\allowbreak{}g\allowbreak{}F\allowbreak{}r\allowbreak{}a\allowbreak{}m\allowbreak{}e\allowbreak{}L\allowbreak{}a\allowbreak{}y\allowbreak{}e\allowbreak{}r\allowbreak{}(\allowbreak{}.\allowbreak{}.\allowbreak{}.\allowbreak{})}
\end{itemize}

\subsection{Why this matters}\label{record-264-why-this-matters}

Since \texttt{v\allowbreak{}4\allowbreak{}.\allowbreak{}0\allowbreak{}.\allowbreak{}0}, the shell has become more explicit about:

\begin{itemize}
\tightlist
\item
  topmost-layer ownership
\item
  dismissal order
\item
  footer/help truth
\item
  settings/search/help/notification/quit layering
\end{itemize}

That means older tests or utilities that inferred shell state from raw render order, footer strings, or branch order inside \texttt{a\allowbreak{}p\allowbreak{}p\allowbreak{}-\allowbreak{}f\allowbreak{}r\allowbreak{}a\allowbreak{}m\allowbreak{}e\allowbreak{}.\allowbreak{}t\allowbreak{}s} are more likely to drift.

\subsection{What to re-test}\label{record-264-what-to-re-test}

\begin{itemize}
\tightlist
\item
  \texttt{E\allowbreak{}s\allowbreak{}c} behavior
\item
  search/help/settings transitions
\item
  notification center and history behavior
\item
  quit confirmation flows
\item
  mouse shielding while exclusive overlays are open
\end{itemize}

\section{5. Runtime Internals: Prefer Engine Seams Over Shell Folklore}\label{record-264-runtime-internals-prefer-engine-seams-over-shell-folklore}

If you extend or inspect runtime internals directly, start moving toward the explicit runtime-engine exports introduced since \texttt{v\allowbreak{}4\allowbreak{}.\allowbreak{}0\allowbreak{}.\allowbreak{}0}.

The relevant surfaces are the \texttt{R\allowbreak{}E} slices:

\begin{itemize}
\tightlist
\item
  state machine and view stack ownership
\item
  retained layouts and invalidation
\item
  layout-driven input routing
\item
  command/effect buffering
\item
  component layout and interaction contracts
\end{itemize}

This does \textbf{not} mean the migration is complete. The repo still tracks \href{https://github.com/flyingrobots/bijou/blob/cbd70612140e7a450d764b3c7be1c6af878b99bc/docs/method/graveyard/legacy-backlog/up-next/RE-007-migrate-framed-shell-onto-runtime-engine-seams.md}{RE-007} because some framed-shell behavior still needs to move more fully onto those seams.

The practical rule is:

\begin{itemize}
\tightlist
\item
  prefer exported runtime seams where they exist
\item
  avoid adding new shell-only branch-order logic if you can use the runtime-engine model instead
\end{itemize}

\section{6. Localization Is Opt-In, Not Forced}\label{record-264-localization-is-opt-in-not-forced}

If you do not need localization, you do not need to migrate code just because the new localization packages exist.

If you do want localization, \texttt{v\allowbreak{}4\allowbreak{}.\allowbreak{}1\allowbreak{}.\allowbreak{}0} gives you a real substrate:

\begin{itemize}
\tightlist
\item
  runtime locale/direction contracts
\item
  tools for catalogs and exchange workflows
\item
  Node filesystem adapters
\item
  XLSX-oriented tooling
\item
  localized shell chrome and DOGFOOD surfaces
\end{itemize}

That is an adoption opportunity, not a forced migration.

\section{7. Docs And DOGFOOD Usually Require No App Migration}\label{record-264-docs-and-dogfood-usually-require-no-app-migration}

Most of the DOGFOOD, design-system, METHOD, and release-doc work does not require code changes in downstream apps.

It does change what is now available to you:

\begin{itemize}
\tightlist
\item
  a much more complete field guide
\item
  stronger component-family guidance
\item
  clearer repo doctrine and release documentation
\end{itemize}

Treat these as improved references, not breaking changes.

\section{8. Validation Checklist After Upgrading}\label{record-264-validation-checklist-after-upgrading}

For a normal app upgrade, run at least:

\begin{CodeBlock}
npm run build
npm test
\end{CodeBlock}

If your app uses TypeScript and the fullscreen runtime, also run your typecheck and any shell smoke tests you have.

Pay special attention to:

\begin{itemize}
\tightlist
\item
  framed shell compile errors
\item
  command typing errors
\item
  tests around help/settings/search/quit
\item
  notification history and overlay behavior
\item
  mouse interaction over modal/help/palette surfaces
\end{itemize}

\section{9. Migration Recommendations By App Shape}\label{record-264-migration-recommendations-by-app-shape}

\subsection{Core CLI / prompt app}\label{record-264-core-cli-prompt-app}

Usually:

\begin{itemize}
\tightlist
\item
  upgrade dependencies
\item
  rebuild
\item
  read the new docs when you want deeper design-system guidance
\end{itemize}

\subsection{Fullscreen TUI app without heavy shell customization}\label{record-264-fullscreen-tui-app-without-heavy-shell-customization}

Usually:

\begin{itemize}
\tightlist
\item
  upgrade dependencies
\item
  fix any framed typing fallout
\item
  re-test shell overlays and interaction flows
\end{itemize}

\subsection{Fullscreen app with custom shell helpers or runtime extensions}\label{record-264-fullscreen-app-with-custom-shell-helpers-or-runtime-extensions}

Usually:

\begin{itemize}
\tightlist
\item
  audit command typings
\item
  audit shell-layer assumptions
\item
  move toward exported framed/runtime helpers
\item
  re-run high-signal interaction tests
\end{itemize}

\subsection{Localization-aware app}\label{record-264-localization-aware-app}

Usually:

\begin{itemize}
\tightlist
\item
  decide whether to adopt the new i18n packages now or later
\item
  verify package availability in the published release
\item
  add localized shell copy incrementally instead of rewriting the whole app at once
\end{itemize}

\section{Final Advice}\label{record-264-final-advice}

Compared with the \texttt{v\allowbreak{}4\allowbreak{}.\allowbreak{}0\allowbreak{}.\allowbreak{}0} migration, \texttt{v\allowbreak{}4\allowbreak{}.\allowbreak{}1\allowbreak{}.\allowbreak{}0} is less about one hard boundary change and more about aligning your app with clearer repo truth.

The important upgrades are:

\begin{itemize}
\tightlist
\item
  accept the typed framed-shell surface
\item
  accept the more honest \texttt{C\allowbreak{}m\allowbreak{}d\allowbreak{}<\allowbreak{}M\allowbreak{}>} contract
\item
  trust exported shell/runtime inspection helpers over folklore
\item
  adopt localization only if and when you need it
\end{itemize}

If your app is mostly core Bijou plus light TUI usage, the migration should be manageable. If you have deep custom shell/runtime code, this is a good release to pay down those assumptions.

